OpenAI側では、2月2日のCodex app、2月5日のGPT-5.3-CodexとSystem Card、2月12日のGPT-5.3-Codex-Sparkという順で更新が続いた。ここで重要なのは、四つを一つの『新モデル』として潰さないことだ。Codex appは複数のcoding agentと長時間の並列作業を管理するworkspace、GPT-5.3-Codexはcodingに加えてresearch・tool use・複雑なexecutionを担うモデル、System Cardはその安全性評価、Sparkはrealtime coding interactionを狙うresearch previewであり、同じ流れに属しながら役割は異なる。\autocite{src-15823ae5556e,src-4a841ea7ec7c,src-9aba7995b6a5,src-dd64fe3084a4}


Codex appが示した変化は、AI codingの操作単位を一往復のprompt-responseから、複数agentが並行して進める作業へ移した点にある。GPT-5.3-Codexも長時間のresearchやtool use、複雑なexecutionを対象に掲げており、UIとモデルの双方で『一度答える』より『仕事を持ち続ける』方向が明確になった。これはcoding benchmarkの更新だけでは説明しにくい、実行系としての設計変更である。\autocite{src-15823ae5556e,src-9aba7995b6a5}


Anthropicも同じ月に、2月5日のClaude Opus 4.6と2月17日のClaude Sonnet 4.6を通じて、codingだけでなくsustained agentic work、computer use、agent planning、knowledge workへ重点を広げた。Sonnet 4.6では1M-token context windowが示されたが、発表時点ではbetaである。Opus 4.6とSonnet 4.6は二つの独立記事に分解するより、2月中に続いた同一方向の更新として読む方が技術的な連続性を捉えやすい。\autocite{src-f6aa679babd7}


両社の一次資料を共通軸で並べると、競争の焦点は『codeを生成できるか』から、長いcontextを保持し、toolやcomputerを使い、途中状態を抱えたまま作業を継続できるかへ移っている。OpenAIでは複数agentを扱うworkspaceとlong-running execution、Anthropicではsustained agentic workとcomputer use・planningが前面に出た。2月をagentic engineeringへの収束として捉える根拠は、この実行時間と作業面の拡張にある。\autocite{src-15823ae5556e,src-9aba7995b6a5,src-f6aa679babd7}


同じagentic engineeringでも、Sparkは別の枝を示す。GPT-5.3-Codex-Sparkはrealtime coding interaction向けのresearch previewとして位置付けられた。OpenAIは1,000 tokens/s超というthroughputを掲げているが、これはCerebrasの専用hardwareとresearch-preview条件に依存するvendor claimである。したがって『通常のCodexが常時その速度で動く』と一般化してはならない。ここで見えるのは、長時間自律実行とは反対側にある、人間との低latencyな往復を最適化する設計軸だ。\autocite{src-15823ae5556e,src-4a841ea7ec7c,src-9aba7995b6a5,src-dd64fe3084a4}


